\documentclass[10pt]{article}

\usepackage[margin=0.78in]{geometry}
\usepackage{amsmath}
\usepackage{booktabs}
\usepackage{enumitem}
\usepackage{hyperref}

\setlist[itemize]{leftmargin=1.2em,itemsep=0.1em,topsep=0.2em}
\setlength{\parindent}{0pt}
\setlength{\parskip}{0.32em}

\title{\vspace{-1.2cm}From Sampling Shocks to Discovery Outcomes:\\Downstream Linkage Options}
\author{}
\date{\vspace{-0.8cm}May 2026}

\begin{document}
\maketitle
\vspace{-1.2cm}

\section*{Three Established Facts}

Current regression results (Tables 1--5, FC3, FC5; 100\,km cell $\times$ year panel, 2005--2023) establish three robust empirical patterns:

\begin{enumerate}[label=(\alph*)]
    \item \textbf{Conflict reduces biodiversity sampling.} The effect survives saturated fixed effects (cell + country$\times$year + biome$\times$year), is robust on both extensive and intensive margins, and cumulates: the sum of contemporaneous + two lags is 2--3$\times$ the contemporaneous coefficient alone.
    \item \textbf{The effect is amplified in biodiverse cells.} The conflict$\times$richness interaction is consistently negative and significant: conflict reduces sampling most where species richness is highest.
    \item \textbf{Conflict selectively deters foreign collectors.} Domestic collecting is essentially unaffected; the aggregate decline is driven by withdrawal of foreign (especially distant/cross-continent) collectors. The binary conflict measure ($\mathbf{1}[\text{events}>0]$) produces $\sim$2$\times$ larger coefficients than log(1+events).
\end{enumerate}

BIN-discovery regressions show that conflict reduces new BIN discovery, but this effect is almost entirely mediated by sampling volume---controlling for $\log(1+\text{records})$ eliminates the conflict coefficient on new BINs. The mechanism is: conflict $\to$ less sampling $\to$ fewer discoveries, not conflict $\to$ lower-quality sampling.

The question is whether we can link these sampling shocks to measurable \emph{downstream} research and discovery outcomes.

\section*{Option A: Publication and Sequence Linkage}

\textbf{Chain:} BOLD specimen $\to$ GenBank accession $\to$ NCBI PubMed citations $\to$ publication counts, journals, time-to-publication.

\textbf{Data linkage.} The raw BOLD TSV files contain the field \texttt{insdc\_acs} (INSDC/GenBank accession number), which was dropped during the minimal-records build (\texttt{00\_build\_bold\_minimal.py}) but can be re-extracted. GenBank accessions provide a clean, non-fuzzy join key to the NCBI ecosystem:
\begin{itemize}
    \item \textbf{Accession $\to$ PubMed:} NCBI Entrez \texttt{elink} API maps nucleotide accessions to citing PubMed IDs in bulk.
    \item \textbf{Sequence metadata:} GenBank records carry deposit date (measures collection-to-sequencing delay), organism, and submitting institution.
    \item \textbf{Citation graph:} PubMed IDs link to publication date, journal, MeSH terms, and forward citations.
\end{itemize}

\textbf{Feasibility: High.} No new raw data downloads beyond NCBI API queries (free, well-documented, bulk-capable).

\textbf{insdc\_acs fill rates (sampled, May 2026 audit).} Coverage is highly heterogeneous across taxa:

\begin{center}
\small
\begin{tabular}{lr}
\toprule
\textbf{Taxon (BOLD)} & \textbf{insdc\_acs fill} \\
\midrule
Fungi & 95\% \\
Chordata & 87\% \\
Mollusca & 73\% \\
Insects (family-level samples, Lep/Hym/Col/Hem/Dip)\hspace{1em} & 17--46\% (median $\approx$ 33\%) \\
Plantae & 28\% \\
\bottomrule
\end{tabular}
\end{center}

Since insects dominate BOLD by record count, the all-animal weighted fill rate is approximately 35--40\%. Option A is therefore primarily a Chordata + Fungi story (cleanly linkable), with insects, mollusks, and plants entering as a noisier secondary layer.

\textbf{Outcomes.} Cell$\times$year panel of: (i) sequences deposited in GenBank, (ii) publications citing those sequences, (iii) time from collection to publication, (iv) journal quality.

\textbf{Limitations.} Many BOLD records lack GenBank accessions, with the gap concentrated in plants and some insect families. Publications citing a sequence are not always ``discoveries''---many are phylogenetic surveys or methods papers. The causal chain is long: conflict $\to$ fewer specimens $\to$ fewer sequences $\to$ fewer papers.

\section*{Option B: Natural Products and Applied Discovery}

\textbf{Chain:} Sampled species/genus $\to$ natural product databases $\to$ known compounds $\to$ bioactivity assays $\to$ patents/clinical leads.

\textbf{Data linkage.} The join key is \emph{species or genus name}, not a shared identifier. Three open natural-product databases provide species$\to$compound mappings:
\begin{itemize}
    \item \textbf{LOTUS} ($\sim$300K compound--taxon pairs). Fully open, Wikidata-linked, cleanest taxonomy. Best starting point.
    \item \textbf{COCONUT} ($\sim$400K compounds, $\sim$50K species links). Downloadable as CSV/SDF. Broader compound coverage.
    \item \textbf{KNApSAcK} ($\sim$50K metabolite--species pairs). Complementary, especially for plant secondary metabolites.
\end{itemize}
Downstream bioactivity and translational potential come from PubChem (compound $\to$ assay results) and ChEMBL (compound $\to$ drug-target activity). Patent linkage is possible via PubChem patent annotations or Google Patents.

\textbf{Construction.}
\begin{enumerate}
    \item Download LOTUS $\to$ extract species $\to$ compound count and compound novelty per species.
    \item Cross-reference with BOLD/GBIF species lists $\to$ which sampled species have known natural products.
    \item Build cell-level ``chemical potential'' from species richness $\times$ average compound yield per species in that cell's species pool.
    \item Test: does conflict reduce sampling of chemically valuable species disproportionately?
\end{enumerate}

\textbf{Feasibility: Medium.} Name matching between BOLD taxonomic names and natural-product databases is fuzzy (synonyms, spelling variants, rank disagreements). Coverage is heavily concentrated in plants and fungi, which aligns well with the GBIF Plantae layer and BOLD fungi but limits applicability to animal taxa. Requires building a species-level ``potential discovery value'' measure, then aggregating to cell$\times$year.

\textbf{Story.} This is the ``option value'' argument: an unsampled species in a conflict zone could harbor undiscovered compounds. The framing---conflict $\to$ lost drug candidates---is a stronger headline than conflict $\to$ fewer papers.

\section*{Comparison}

\begin{center}
\begin{tabular}{lll}
\toprule
 & \textbf{A: Publication linkage} & \textbf{B: Natural products} \\
\midrule
Construction effort & Low (re-extract 1 field + API) & Medium (name matching + multiple DBs) \\
Join quality & Clean (GenBank accession) & Fuzzy (species names) \\
Time to first result & 1--2 weeks & 3--4 weeks \\
Narrative strength & ``Conflict $\to$ fewer papers'' & ``Conflict $\to$ lost drug candidates'' \\
Novelty & Moderate & High \\
Strongest taxa & Chordata, Fungi (high \texttt{insdc\_acs}) & Plantae, Fungi (NP-DB coverage) \\
Weakest taxa & Plantae (28\% \texttt{insdc\_acs}) & Animals (sparse NP entries) \\
\bottomrule
\end{tabular}
\end{center}

The two paths are taxonomically \emph{complementary}, not redundant. Option~A is mainly an animal + fungi story anchored on BOLD's GenBank deposition; Option~B is mainly a plant + fungi story anchored on natural-product databases. \textbf{Fungi sit at the intersection of both layers}: 95\% \texttt{insdc\_acs} coverage in BOLD plus strong representation in COCONUT/LOTUS/KNApSAcK make fungi the natural cross-validation taxon---one taxon where both downstream pipelines work, allowing internal consistency checks. A reasonable sequencing is to start with~A on Chordata + Fungi (clean linkage, fast first result) while building~B in parallel on Plantae + Fungi.

\end{document}
